% ECONOMICS_PROBLEM: {"id": "I11-561", "title": "Hospital and physician consolidation", "jel_code": "I11", "source_id": "OP-0054"}
% !TeX program = xelatex
\documentclass[11pt,a4paper]{article}
\usepackage[margin=23mm]{geometry}
\usepackage{fontspec}
\setmainfont{Latin Modern Roman}
\usepackage{amsmath,amssymb,mathtools}
\usepackage{xcolor,enumitem,booktabs,longtable,array,xurl,hyperref}
\definecolor{muted}{HTML}{53636C}
\hypersetup{colorlinks=true,urlcolor=blue,linkcolor=blue}
\setlength{\parindent}{0pt}
\setlength{\parskip}{5pt}
\newcommand{\R}{\mathbb R}
\newcommand{\E}{\mathbb E}
\newcommand{\N}{\mathbb N}
\newcommand{\ind}{\mathbf 1}
\newcommand{\Var}{\operatorname{Var}}
\newcommand{\Cov}{\operatorname{Cov}}
\newcommand{\supp}{\operatorname{supp}}
\DeclareMathOperator*{\argmin}{arg\,min}
\DeclareMathOperator*{\argmax}{arg\,max}
\newcommand{\entrymeta}[3]{{\sffamily\small\color{muted}\textbf{\texttt{#1}}\quad|\quad #2\par #3\par}}
\newcommand{\Problemref}[1]{\texttt{#1}}
\newcommand{\JELref}[2]{\texttt{#2} (JEL \texttt{#1})}
\begin{document}
\section*{Hospital and physician consolidation}
\textbf{Problem ID:} \texttt{I11-561}\quad \textbf{JEL:} \texttt{I11}\par
% BEGIN ECONOMICS STATEMENT
\label{problem:OP-0054}
\label{atlas:prob:H4}
\noindent\textbf{Original atlas identifier:} H4; entry 54. \textbf{Field:} Health economics. \textbf{Original tractability label:} Medium.

\noindent\textbf{Classification:} Parameterized theoretical/empirical research agenda; not a certified unsolved theorem.

\subsection*{Definitions and assumptions}
Let $M\in\{0,1\}$ indicate a specified hospital or physician merger relative to a baseline ownership structure. A model $m$ determines insurer--provider bargaining, network inclusion, treatment, capacity, staffing, and quality after each ownership regime; select an equilibrium explicitly. Let $EV_i(M;m)$ be patients' money-metric welfare, including premiums, health quality, travel, and access, excluding the profit changes separately allocated to owners. Let $\Pi_h(M;m)$ be provider and insurer profits net of real costs, assigned to their ultimate owners rather than counted both as firm and household income. Fix weights $w_i,w_h\geq0$ and a welfare boundary covering affected rivals and patients. Let $G(M;m)$ be net public expenditure with monetary social cost $\lambda G$, under a specified public financing rule.

\subsection*{Formal research question}
\emph{Editorial formulation: the mathematical target below is not a theorem quoted from the references.}

Identify the merger-specific welfare contrast
\[
 \Delta W(m)=\sum_iw_i\Delta EV_i(m)+\sum_hw_h\Delta\Pi_h(m)-\lambda\Delta G(m),
\]
where $\Delta$ denotes regime $M=1$ minus $M=0$. Construct a prospective prediction set using premerger data and declared restrictions on bargaining, quality choice, and cost efficiencies; test its coverage against ex-post outcomes. Establish which counterfactual ownership comparison is identified when mergers are endogenous, allowing rival effects and market redefinition. Decompose price changes into transfers and associated utilization/access distortions; define “efficiency” as a real-resource or quality gain, not simply a merger-induced profit increase. If clinical quality is incompletely observed, report welfare bounds rather than only a price effect.

\subsection*{Interpretation and scope}
A consolidation event is not a single homogeneous treatment: horizontal mergers, vertical integration, and physician acquisitions change different choice sets and bargaining positions. The formal target therefore specifies the assets and institutions affected. Postmerger price increases can coexist with cost savings, but their welfare trade-off depends on premiums, coverage, treatment, distributional weights, and financing. Quality should use clinical risk-adjusted outcomes and access measures, with the risk-adjustment assumptions stated. Rival hospitals can change prices and staffing, invalidating designs that assume unaffected comparison units without justification. Prospective validation requires freezing the prediction before observing postmerger outcomes. The intended extension is a merger-specific, clinically informed welfare assessment and calibrated forecast, not a request to decide whether every merger is good or bad.

\subsection*{Source audit and status}
[S1] treats merger endogeneity, [S2] reviews health-care industrial organization, and [S3] studies privately insured prices. [S3] has online publication in 2018 and journal issue year 2019. These antecedents support mechanisms and specific estimates, not a universal present-day consolidation effect or a verified open theorem.

\subsection*{References}
\begin{enumerate}[label={[S\arabic*]},leftmargin=*,itemsep=0.6em]
\item Leemore Dafny (2009). Estimation and Identification of Merger Effects: An Application to Hospital Mergers. Journal of Law and Economics 52(3), 523--550. DOI: 10.1086/600079.
\url{https://www.journals.uchicago.edu/doi/10.1086/600079}
\newline\emph{Locator and support limits:} Abstract: rival analysis and merger endogeneity; identifying assumptions are setting-specific.
\item Martin Gaynor, Kate Ho, and Robert J. Town (2015). The Industrial Organization of Health-Care Markets. Journal of Economic Literature 53(2), 235--284. DOI: 10.1257/jel.53.2.235.
\url{https://www.aeaweb.org/articles?id=10.1257/jel.53.2.235}
\newline\emph{Locator and support limits:} Abstract and review: competition, price, quality, and treatment; a review, not a general merger theorem.
\item Zack Cooper, Stuart V. Craig, Martin Gaynor, and John Van Reenen (2019). The Price Ain’t Right? Hospital Prices and Health Spending on the Privately Insured. The Quarterly Journal of Economics 134(1), 51--107. DOI: 10.1093/qje/qjy020.
\url{https://academic.oup.com/qje/article-abstract/134/1/51/5090426}
\newline\emph{Locator and support limits:} Abstract: private-insurance prices and mergers. Online publication 4 September 2018; volume year 2019.
\end{enumerate}

\subsection*{Common conventions: Source mathematical conventions}

Unless an entry states otherwise, agent and firm sets are finite. State and action spaces are measurable, strategies and outcome maps are measurable, probability laws are specified, and expectations exist for the targets under discussion. Infinite-horizon discount factors lie in $(0,1)$ and discounted objectives are finite. Norms, tolerances and complexity input sizes must be specified for computational comparisons. Constants or primitives that appear locally are inputs, not empirical facts asserted by this catalogue.

\textbf{Identification.} A statistical model $m\in\mathcal M$ implies an observable law $P_m$ and target $\theta(m)$. At law $P$, its identified set is
\[
\Theta(P;\mathcal M)=\{\theta(m):m\in\mathcal M,\ P_m=P\}.
\]
Point identification holds when this set is a singleton, or equivalently when $P_{m_1}=P_{m_2}$ implies $\theta(m_1)=\theta(m_2)$ for all compatible models. Sharp bounds mean bounds attained, or approached when the boundary is not attained, by compatible models. Writing this set is a definition, not a solution: the substantive task is to establish informative restrictions and derive or estimate the set. An unrestricted-confounding model may give only trivial bounds.

\textbf{Inference.} Identification, estimability and valid uncertainty quantification are separate questions. Fix $\alpha\in(0,1)$. A confidence procedure $C_n$ over a declared law class $\mathcal P$ has asymptotic uniform coverage at level $1-\alpha$ when
\[
\liminf_{n\to\infty}\inf_{P\in\mathcal P}\Pr_P\{\theta(P)\in C_n\}\geq1-\alpha.
\]
For set-identified targets the coverage event must be defined separately, for example coverage of the full identified set. If an entry uses identification-indexed models instead of uniquely indexed laws, the same coverage convention is applied over those models. Pointwise limits do not establish uniform validity, and asymptotic validity does not guarantee finite-sample precision. Sample sizes, dependence structures and weak-identification sequences are part of each application.

\textbf{Counterfactuals.} $Y_i(a)$ is a potential outcome under a specified intervention, including its timing, implementation and versions. It is not $Y_i$ conditional on voluntarily choosing $a$. A full intervention vector permits interference; shorthand scalar treatment notation does not silently impose no interference. Assignment, selection, attrition, anticipation and measurement must be specified. A claim about an instrument's local effect is not automatically a population policy effect.

\textbf{Economic equilibrium.} A counterfactual model includes best responses, constraints, feasibility, market clearing, state transitions and expectations consistent with the stated information. When several equilibria exist, either a declared selector is an input or the target is set-valued. Comparisons across policy regimes must use the stated selection convention. Reduced-form relationships are not presumed invariant after an equilibrium-changing intervention.

\textbf{Optimization and welfare.} Optimization is over the stated feasible set. If attainment is not established, $\sup$ or $\inf$ and $\varepsilon$-optimal choices replace an asserted maximizer or minimizer. For example, with value $V=\sup_{a\in\mathcal A}W(a)<\infty$, an $\varepsilon$-optimal set is $\{a\in\mathcal A:W(a)\geq V-\varepsilon\}$ for $\varepsilon>0$. Benefits, costs, risk and health or environmental outcomes can be aggregated only using declared units and conversion weights. Interpersonal utility weights, the treatment of fiscal transfers and foreign profits, discounting, and the behavioral welfare criterion are explicit inputs. A comparative-static derivative additionally requires existence and appropriate differentiability of the selected counterfactual.

\textbf{Sources.} Local labels [S1], [S2], and so forth restart in every question. Locators identify the relevant abstract, model, section or result at the level actually checked; a bibliographic or abstract-level verification is not represented as inspection of an entire proof. Working papers are distinguished from later publications and changing versions. Source summaries are paraphrased. All URL checks and editorial formulations refer to the audit date stated above.
% END ECONOMICS STATEMENT
\end{document}
